% Part:sets-functions-relations
% Chapter: size-of-sets
% Section: pairing-alt

\documentclass[../../../include/open-logic-section]{subfiles}

\begin{document}

\olfileid{sfr}{siz}{pai-alt}

\olsection{एक पर्यायी जोडीकरण फलन}

\begin{explain}
$\Nat^2$ ची अशी इतर प्रगणनेही आहेत की त्यांची व्युत्क्रमे शोधणे अधिक
सोपे होते. त्यांपैकी एक पुढे दिले आहे. प्रगणन सरणीत दाखवण्याऐवजी,
सुरुवातीला रिकाम्या असलेल्या जागांशी जोडलेल्या धन पूर्णांकांच्या
यादीपासून सुरुवात करा. या जागा पुढीलप्रमाणे एकामागून एक
$\tuple{n,m}$ जोड्यांनी भरत आहोत, अशी कल्पना करा. पहिल्या स्थानी~$0$
असलेल्या जोड्यांपासून (म्हणजे $\tuple{0,m}$ जोड्यांपासून) सुरुवात करा.
पहिली जोडी (म्हणजे $\tuple{0,0}$) पहिल्या रिकाम्या जागेत ठेवा; मग एक
रिकामी जागा सोडून दुसरी जोडी पुढील रिकाम्या जागेत ठेवा आणि पुन्हा एक
जागा सोडा; ही प्रक्रिया अशीच पुढे चालू ठेवा. गोठवलेल्या इंग्रजी
मजकुरात दुसरी जोडी $\tuple{0,2}$ अशी दिली आहे; परंतु खालील सारणीनुसार
येथे 〈0,1〉 अभिप्रेत आहे. आता आपल्या प्रगणनाची (अपूर्ण) सुरुवात अशी
दिसते:
%READERNOTE{MRSIZ-003: गोठवलेल्या इंग्रजी मजकुरात एक रिकामी जागा सोडल्यानंतरची दुसरी जोडी 〈0,2〉 अशी दिली आहे, पण लगेचची सारणी आणि वर्णनाची क्रमवार रचना ती 〈0,1〉 असावी हे दाखवतात. सूत्र-संरचना जशीच्या तशी ठेवून मराठी मजकुरात अपेक्षित जोडी स्पष्ट केली आहे.}
\[\small
\begin{array}{@{}c c c c c c c c c c c@{}}
\mathbf 1 & \mathbf 2 & \mathbf 3 & \mathbf 4 & \mathbf 5 & \mathbf 6 & \mathbf 7 & \mathbf 8 & \mathbf 9 & \mathbf{10} & \dots \\ \\
\tuple{0,0} &  & \tuple{0,1} &  & \tuple{0,2} &  & \tuple{0,3} & & \tuple{0,4} &  & \dots \\
\end{array}
\]
अजून रिकाम्या राहिलेल्या जागांसाठी $\tuple{1,m}$ जोड्यांबाबत हीच
प्रक्रिया पुन्हा करा; पुन्हा प्रत्येक दुसरी रिकामी जागा वगळा:
\[\small
\begin{array}{@{}c c c c c c c c c c c@{}}
\mathbf 1 & \mathbf 2 & \mathbf 3 & \mathbf 4 & \mathbf 5 & \mathbf 6 & \mathbf 7 & \mathbf 8 & \mathbf 9 & \mathbf{10} & \dots \\ \\
\tuple{0,0} & \tuple{1,0} & \tuple{0,1} &  & \tuple{0,2} & \tuple{1,1} & 
\tuple{0,3} & & \tuple{0,4} &  \tuple{1,2} & \dots \\
\end{array}
\]
यानंतर $\tuple{2,m}$ जोड्या आणि मग $\tuple{3,m}$ जोड्या त्याच पद्धतीने
भरा; त्यापुढेही पहिला घटक अनुक्रमे 4, 5, इत्यादी करत पुढे जायचे आहे.
आपल्या पूर्ण प्रगणनाची सुरुवात म्हणून अशी दिसते:
%READERNOTE{OLSIZ-003: गोठवलेल्या इंग्रजी मजकुरात “pairs 〈2,m〉, 〈2,m〉, etc.” अशी पुनरुक्ती आहे. पूर्ण सारणीतील 〈3,0〉 आणि पुढील सर्व ओळींवरून दुसऱ्या ठिकाणी 〈3,m〉 अभिप्रेत असल्याचे निश्चित होते; मराठी मजकुरात दुसरी जोडी त्यानुसार दुरुस्त केली आहे.}
\[\small
\begin{array}{@{}cc c c c c c c c c c@{}}
\mathbf 1 & \mathbf 2 & \mathbf 3 & \mathbf 4 & \mathbf 5 & \mathbf 6 & \mathbf 7 & \mathbf 8 & \mathbf 9 & \mathbf{10} & \dots \\ \\
\tuple{0,0} & \tuple{1,0} & \tuple{0,1} & \tuple{2,0}  & \tuple{0,2} & 
\tuple{1,1} & \tuple{0,3} & \tuple{3,0}  & \tuple{0,4} &  \tuple{1,2} & \dots \\
\end{array}
\]
वरील सरणीतील चौकटींना या प्रगणनानुसार क्रमांक दिले, तर आपल्याला
नीट नागमोडी रेषा दिसणार नाही; त्याऐवजी पुढील मांडणी दिसेल:
\[
\begin{array}{ c | c | c | c | c | c | c | c }
& \mathbf 0 & \mathbf 1 & \mathbf 2 & \mathbf 3 & \mathbf 4 & \mathbf 5 & \dots \\
\hline
\mathbf 0 & 1 & 3 & 5 & 7 & 9 & 11 & \dots \\
\hline
\mathbf 1 & 2 & 6 & 10 & 14 & 18 & \dots & \dots \\
\hline
\mathbf 2 & 4 & 12 & 20 & 28 & \dots & \dots & \dots \\
\hline
\mathbf 3 & 8 & 24 & 40 & \dots & \dots & \dots & \dots \\
\hline
\mathbf 4 & 16 & 48 & \dots & \dots & \dots & \dots & \dots \\
\hline
\mathbf 5 & 32 & \dots & \dots & \dots & \dots & \dots & \dots \\
\hline
\vdots & \vdots & \vdots & \vdots & \vdots & \vdots & \vdots & \ddots\\
\end{array}
\]
आपल्या प्रगणनात ओळ~$0$ मधील जोड्या विषम क्रमांकांच्या स्थानी आहेत,
हे दिसते; म्हणजे $\tuple{0,m}$ ही जोडी $2m+1$ या स्थानी आहे. दुसऱ्या
ओळीतील $\tuple{1,m}$ जोड्या विषम संख्येच्या दुप्पट क्रमांक असलेल्या
स्थानी आहेत; विशेषतः त्या $2 \cdot (2m+1)$ या स्थानी आहेत. तिसऱ्या
ओळीतील $\tuple{2,m}$ जोड्या विषम संख्येच्या चौपट क्रमांक असलेल्या
स्थानी, म्हणजे $4 \cdot (2m+1)$ या स्थानी आहेत; आणि ही रचना अशीच
पुढे चालते. प्रत्येक ओळीत $(2m+1)$ च्या गुणकांचे मूल्य $1$, $2$,
$4$, $8$, \dots{} असे आहे; हे नेमके~$2$ चे घात आहेत:
$1= 2^0$, $2 = 2^1$, $4 = 2^2$, $8 = 2^3$, \dots\@ विचाराधीन
जोडीचा पहिला घटक हाच नेहमी संबंधित घातांक असतो. म्हणून
$\tuple{n,m}$ या जोडीसाठी गुणक $2^n$ आहे. यातून आपल्याला
$2^n \cdot (2m+1)$ हे सर्वसाधारण सूत्र मिळते. पण या फलनाने जोड्या
\emph{धन} पूर्णांकांकडे पाठवल्या जातात; म्हणजे $\tuple{0,0}$ चे
स्थान~$1$ आहे. स्थानांची सुरुवात~$0$ पासून करायची असेल, तर उत्तरातून
$1$ वजा करावा लागतो. त्यामुळे पुढील फलन मिळते:
\end{explain}

\begin{ex}
$h\colon \Nat^2 \to \Nat$ हे पुढीलप्रमाणे दिलेले फलन घ्या:
\[
h(n,m) = 2^n (2m+1) - 1
\]
हे नैसर्गिक संख्यांच्या जोड्यांचा संच~$\Nat^2$ यासाठीचे जोडीकरण फलन
आहे.
\end{ex}

\begin{explain}
त्यानुसार, $\Nat^2$ च्या आपल्या दुसऱ्या प्रगणनात $\tuple{0,0}$ या
जोडीचा संकेतांक $h(0,0) = 2^0(2\cdot 0+1) - 1 = 0$ आहे;
$\tuple{1,2}$ हिचा संकेतांक $2^{1} \cdot (2 \cdot 2 + 1) - 1 = 2
\cdot 5 - 1 = 9$ आहे; आणि $\tuple{2,6}$ हिचा संकेतांक $2^{2} \cdot (2
\cdot 6 + 1) - 1 = 51$ आहे.
\end{explain}

कधी कधी नैसर्गिक संख्यांच्या जोड्यांचे~$\Nat^2$ सांकेतीकरण करताना ते
आच्छादक असण्याची आवश्यकता नसते. अशा सांकेतीकरणांची व्युत्क्रमे केवळ
अंशतः फलने असतात. 

\begin{ex}
$j\colon \Nat^2 \to \Nat^+$ हे पुढीलप्रमाणे दिलेले फलन घ्या:
\[
j(n,m) = 2^n3^m
\]
हे $\Nat^2 \to \Nat$ असे !!a{injective} फलन आहे.
\end{ex}

\end{document}
